% Drawing 80/01-5 -- folder 80, batch 01, archive page 12 (PDF page 13).
% Transcription: en haut de la page, au crayon : cinq pseudo-droites en étoile, avec au centre un pentagone marqué $5$, entouré de cinq triangles marqués $3$ ; autour, cinq régions marquées $4_1$, $4_2$, $4_3$, $4_5$, $4_6$, les mêmes indices étant répétés aux extrémités opposées ; un court trait épais traverse le pentagone d'un sommet à l'autre ; au-dessus, à gauche, deux traits isolés
% Kind: arrangement of five pseudo-lines (pentagram). Method 2 (schematic redraw from the facsimile). Status: ok.
% Notes: Pencil. Five pseudo-lines drawn as smooth curves through the ten crossing points, each point shared by the two curves that cross there. Central pentagon 5, five triangles 3, the regions 4_1, 4_2, 4_3, 4_5, 4_6 and the same labels repeated at the opposite ends (the same regions across infinity). The short heavy stroke across the pentagon and the two isolated strokes at the top left are drawn. Disagreement with the transcription's note: it counts six pseudo-lines and regions 4_1 to 4_6; the page has five long lines (the sixth stroke is the short one inside the pentagon) and no label 4_4: the label at the right end, read 4_6, is the region 4_6 seen across infinity. Redrawn 2026-11-12 (issue 128): to check again.
% Coordinates: 1 unit = 1 mm = 1/100 of the crop width (bbox in page fractions: [0.22, 0.17, 0.6, 0.49]).
% Redrawn by Claude (claude-fable-5-2), issue #128, 2026-11-12; not checked by a human; not compiled (no TeX in the session).
% Rights: reproduces part of the work -- behind RENDER_TRANSCRIPTS, never published by itself.
\documentclass[tikz,border=3pt]{standalone}
\usepackage{amsmath,amssymb}
\usetikzlibrary{patterns,arrows.meta}
\tikzset{pencil/.style={black!45}, struck/.style={black!35}, mpt/.style={circle,draw,inner sep=1.3pt}, pt/.style={circle,fill,inner sep=0.9pt}, opt/.style={circle,draw,fill=white,inner sep=1.1pt}}
\begin{document}
\begin{tikzpicture}[x=1mm,y=1mm,line cap=round,line join=round,font=\scriptsize]
\draw plot[smooth, tension=0.55] coordinates {(10,-30) (20,-35) (30,-40) (40,-45) (49,-50) (56.8,-55) (68,-59) (75,-63) (82,-67) (90,-70)};
\draw plot[smooth, tension=0.55] coordinates {(13,-17) (20,-27) (30,-40) (36,-51) (43,-60) (49.5,-66.8) (56.8,-78.5) (63,-92) (67,-100) (69.5,-106.5)};
\draw plot[smooth, tension=0.55] coordinates {(68,-15) (62,-26) (56,-37) (49,-50) (43,-60) (36,-68) (30.5,-76) (25,-87) (19,-98)};
\draw plot[smooth, tension=0.55] coordinates {(56,-17) (56,-37) (56.8,-55) (57,-63.4) (56.8,-78.5) (57,-99) (58,-112)};
\draw plot[smooth, tension=0.55] coordinates {(8,-90) (20,-82) (30.5,-76) (40,-71) (49.5,-66.8) (57,-63.4) (68,-59) (80,-55.5) (93,-53.5)};
\draw[thick] (49.5,-66.8) -- (56.5,-56.5);
\draw[thick] (13,-6) -- (12.5,-8.5);
\draw (19.5,-5.5) -- (17.5,-13);
\node[inner sep=1pt] at (14,-24) {\small $4_2$};
\node[inner sep=1pt] at (63,-15) {\small $4_1$};
\node[inner sep=1pt] at (46,-31) {\small $4_5$};
\node[inner sep=1pt] at (74,-41) {\small $4_3$};
\node[inner sep=1pt] at (22,-61) {\small $4_6$};
\node[inner sep=1pt] at (41,-52) {\small $3$};
\node[inner sep=1pt] at (53,-48) {\small $3$};
\node[inner sep=1pt] at (50,-58.5) {\small $5$};
\node[inner sep=1pt] at (60,-59.5) {\small $3$};
\node[inner sep=1pt] at (40,-67) {\small $3$};
\node[inner sep=1pt] at (52,-70.5) {\small $3$};
\node[inner sep=1pt] at (89,-62) {\small $4_6$};
\node[inner sep=1pt] at (74,-81) {\small $4_2$};
\node[inner sep=1pt] at (40,-92) {\small $4_1$};
\node[inner sep=1pt] at (8,-98) {\small $4_3$};
\node[inner sep=1pt] at (63,-107) {\small $4_5$};
\end{tikzpicture}
\end{document}
